%!TeX root = ./../Presentation.tex

\section{Ergebnisse}

\begin{frame}[fragile]
    \begin{center}
        \huge\textbf{\insertsection}
    \end{center}
\end{frame}

\begin{frame}[allowframebreaks]
\frametitle{\insertsection}

\begin{block}{QEMU Erweiterung}
\begin{itemize}
    \item Software benötigt nur geringfügige Anpassungen um sowohl in QEMU als
        auch auf Hardware lauffähig zu sein
    \item Nutzung bereits integrierter, kompatibler Peripherie Module möglich
    \item Test der GPIO Zustände aufgrund fehlender Simulation nur mittels
        Tracing möglich
\end{itemize}
\end{block}
\newpage
\begin{block}{EDI Erweiterung}
\begin{itemize}
    \item Firmware benötigt zwingend eine extra Abstraktionsebene und ist nicht
        ohne größere Anpassungen auf realer Hardware lauffähig
    \item Portierung des Demo Projekts zu umständlich und komplex
    \item Peripherie Simulation ist theoretisch einfach erweiterbar
\end{itemize}
\end{block}

\end{frame}

\begin{frame}[fragile]
\frametitle{\insertsection}

\textbf{Ausblick:}
\begin{itemize}
    \item Vervollständigung der bereits vorhandenen Peripherie (GPIO)
    \item Integration komplexer Peripherie wie Ethernet oder CAN
    \item Visualisierung des Mikrocontrollers und Testboards (LEDs, Pins)
\end{itemize}

\end{frame}
